\origpage{95}{N. Robertson, P. Seymour / Journal of Combinatorial Theory, Series B 89 (2003) 77--108}
A partial $\Omega$-patchwork $(G, \mu, \Delta, \phi)$ has \textit{tree-width} $\le w$, where $w \ge 0$, if there is a tree-decomposition $(T, \tau)$ of $G^-$ of width $\le w$ such that $\bar{\pi}(G) \subseteq V(\tau(t))$ for some $t \in V(T)$. If $(P, \mathscr{L})$ is a partial $\Omega$-place, and $P|A$ has tree-width $\le w$ for all $A \in \mathscr{L}$, we say that $P$ is an \textit{enlargement} of $P|\mathscr{L}$ by tree-width $\le w$.

\textbf{5.5.} \textit{Let $P = (G, \mu, \Delta, \phi)$ be an $\Omega$-patchwork, let $w \ge 0$, and let $\mathscr{L}, \mathscr{L}^*$ be rooted locations in $G$, such that $\mathscr{L}^{*-}$ is an enlargement of $\mathscr{L}^-$ by tree-width $\le w$. Then $P | \mathscr{L}^*$ is an enlargement of $P | \mathscr{L}$ by tree-width $\le w$.}

\textbf{Proof.} Let $\mathscr{L} = \{C_1, \dots, C_k\}$ where $C_1, \dots, C_k$ are distinct, and for $1 \le i \le k$ let $D_i = (G \setminus C_i)^-$. Then $\mathscr{L}^- = \{(C_1^-, D_1), \dots, (C_k^-, D_k)\}$. (However, $(C_1^-, D_1), \dots, (C_k^-, D_k)$ may not all be distinct.) Let $P | \mathscr{L}^* = (G^*, \mu^*, \Delta^*, \phi^*) \ (= P^* \text{ say})$, using new elements $e(A)$ ($A \in \mathscr{L}^*$). Choose $\mathscr{L}_1, \dots, \mathscr{L}_k \subseteq \mathscr{L}^*$, mutually disjoint and with union $\mathscr{L}^*$, such that for $1 \le i \le k$, every $(A, B) \in \mathscr{L}_i^-$ satisfies $A \subseteq C_i^-$ and $D_i \subseteq B$, and $\mathscr{L}_i^- \cup \{(D_i, C_i^-)\}$ is a location in $G^-$ of tree-width $\le w$. We claim that for $1 \le i \le k$ and all $A \in \mathscr{L}_i$, $A$ is complemented in $C_i$. For certainly $A^- \subseteq C_i^-$ and $D_i \subseteq (G \setminus A)^-$ since $\mathscr{L}_i^- \cup \{(D_i, C_i^-)\}$ is a location in $G^-$. Moreover,
\begin{equation*}
\bar{\pi}(C_i) \cap V(A) \subseteq V(D_i) \cap V(A) \subseteq V(G \setminus A) \cap V(A) = \bar{\pi}(A);
\end{equation*}
and if $v \in V(A)$ is an end of some $e \in E(C_i) \setminus E(A)$, then $e \in E(G) \setminus E(A)$ and so $v \in \bar{\pi}(A)$. This proves that $A$ is complemented in $C_i$. Consequently, for $1 \le i \le k$, $\mathscr{L}_i$ is a rooted location in $C_i$, and so $(P | C_i, \mathscr{L}_i)$ is an $\Omega$-place; let $P_i = (G_i, \mu_i, \Delta_i, \phi_i)$ be its heart (with ``new'' elements $e(A)$ ($A \in \mathscr{L}_i$), some of the new elements of $P^*$).

\hspace*{12pt}(1) \textit{$P_i$ has tree-width $\le w$.}

\textbf{Subproof.} Let $H$ be the heart of $\mathscr{L}_i^- \cup \{(D_i, C_i^-)\}$. Then $H$ is the hypergraph obtained from
\begin{equation*}
M(G^-, \mathscr{L}_i^- \cup \{(D_i, C_i^-)\}) = C_i^- \cap \bigcap ((G \setminus A)^- \colon A \in \mathscr{L}_i)
\end{equation*}
by adding a new edge with set of ends $V(A \cap B)$ for each $(A, B) \in \mathscr{L}_i^-$, and adding one further new edge with set of ends $V(D_i \cap C_i^-)$ (unless $(D_i, C_i^-) \in \mathscr{L}_i$). Also, $G_i^-$ is obtained from
\begin{equation*}
M(C_i, \mathscr{L}_i) = C_i^- \cap \bigcap ((C_i \setminus A)^- \colon A \in \mathscr{L}_i)
\end{equation*}
by adding a new edge with set of ends $\bar{\pi}(A)$ for each $A \in \mathscr{L}_i$. But
\begin{align*}
C_i^- \cap \bigcap ((G \setminus A)^- \colon A \in \mathscr{L}_i) &= C_i^- \cap \bigcap (C_i^- \cap (G \setminus A)^- \colon A \in \mathscr{L}_i) \\
&= C_i^- \cap \bigcap ((C_i \setminus A)^- \colon A \in \mathscr{L}_i);
\end{align*}
and there is a surjection from $\mathscr{L}_i$ onto $\mathscr{L}_i^-$ such that if $A \in \mathscr{L}_i$ is mapped to $(A', B') \in \mathscr{L}_i^-$ then $A' = A^-$ and $V(A' \cap B') = \bar{\pi}(A)$. Consequently, a hypergraph isomorphic to $G_i^-$ may be obtained from $H$ by deleting an edge with set of ends $\bar{\pi}(C_i) = V(D_i \cap C_i^-)$ (unless $(D_i, C_i^-) \in \mathscr{L}_i$) and adding some new edges, each with the same ends as some edge of $H$. (The latter arise when distinct members of $\mathscr{L}_i$
